\documentclass[11pt]{article}
\usepackage[margin=2.4cm]{geometry}
\usepackage{fontspec}
\setmainfont{DejaVu Sans}
\usepackage{juliamono}
\usepackage{fvextra}
\usepackage{xcolor}
\usepackage{xurl}
\usepackage{hyperref}
\usepackage{amssymb}
\newcommand{\checkbox}{$\square$}
\newcommand{\answerchoices}{\checkbox\ YES \quad \checkbox\ NO \quad \checkbox\ CANNOT DETERMINE}
\newcommand{\decisionexplanation}{%
  \par\smallskip\textbf{If NO, give the reason for the mismatch (required). If
  CANNOT DETERMINE, state what context or evidence is missing:}\par
  \rule{\linewidth}{0.4pt}\par\vspace{7pt}\rule{\linewidth}{0.4pt}\par}
\setlength{\parindent}{0pt}\setlength{\parskip}{5pt}
\setlength{\emergencystretch}{2em}
\DefineVerbatimEnvironment{Vsrc}{Verbatim}{fontsize=\small,frame=single,
  breaklines=true,breakanywhere=true}
\DefineVerbatimEnvironment{Vlean}{Verbatim}{fontsize=\scriptsize,frame=single,
  numbers=left,numbersep=5pt,breaklines=true,breakanywhere=true}
\hypersetup{colorlinks=true,urlcolor=blue}

\begin{document}
\begin{center}{\LARGE Proof Verification Document}\\[2pt]{\large reddit\_learnmath\_14s0k9c}\\[2pt]{\small Reddit proof-check posts (newly confirmed)}\end{center}
\begin{center}\fcolorbox{blue!60!black}{blue!4}{\parbox{0.92\linewidth}{\textbf{Make three independent decisions.} Follow the evidence chain: original source $\longrightarrow$ challenged claim $\longrightarrow$ proposed counterexample $\longrightarrow$ Lean theorem. A record is valid only when all three links are correct. Complete the human sign-off before reading the automated review appendix; those earlier results are background, not an answer key.}}\end{center}
\textbf{Use only this document, except for the public source link.} Everything needed for Questions 2 and 3 is included here. For Question 1, open the public source link when one is shown. If there is no public link, the complete stored problem and solution are included here. If the information is still not enough, choose CANNOT DETERMINE.
For each question, choose \textbf{YES}, \textbf{NO}, or \textbf{CANNOT DETERMINE}. Choose CANNOT DETERMINE if the document does not provide enough information for you to decide; do not mark YES if you are uncertain. If you find a specific mismatch, choose NO and explain it in the space below that question. Do not rewrite the claim, add assumptions, or substitute a different counterexample.
\textbf{How copied text is shown.} The problem, quoted passage, and proposed counterexample come from the stored record. The Lean theorem statement and Lean files are copied exactly from the certificate. Some long lines continue on the next printed line; this is only page formatting and does not change the copied text. Line numbers are added only to help you find passages. If the stored record contains a character that cannot be printed, the document shows \texttt{[U+XXXX CONTROL]} in its place.
\begin{description}
\item[Challenged claim.] The exact statement from the source that this record says is false.
\item[Proposed counterexample.] A concrete case that should satisfy the claim's assumptions but make its conclusion false.
\item[Lean certificate.] A Lean file proving a formal theorem. Successful compilation proves the theorem as written, not that it is the right translation of the source.
\end{description}
\section*{Question 1 --- Does this proof verification document accurately represent the source?}
\textbf{What to decide.} Compare the challenged passage below with the original source---either the public link or the complete source text in this document. Answer YES only if they make the same claim and include the same assumptions or conditions. If this document changes the meaning, leaves out a condition, or points to a different passage, answer NO.
\subsection*{Original source}
\textbf{Source record:}
\begin{Vsrc}
qedbench_certificate_availability_triage/data/records_reddit_exhaustive.jsonl
\end{Vsrc}
\textbf{Public source:}\par
\url{https://www.reddit.com/r/learnmath/comments/14s0k9c/checking_for_error_in_group_theory_proof/}\par
\subsection*{Problem or question from the source (verbatim)}
\begin{Vsrc}
checking for error in group theory proof
\end{Vsrc}
\subsection*{Challenged passage recorded from the source (stored text)}
\begin{Vsrc}
So x and y are both mapped to transpositions (a b) and (c d).
\end{Vsrc}
\subsection*{Surrounding source lines (stored text with added line labels)}
The labels refer to lines in the solution text stored with the source record; the source-line content is not transliterated.
\begin{Vsrc}
0001 | Hi everyone, I was recently doing some algebra problems and came up with this solution that looks a bit different from the ones I saw after checking for people's solutions.  Was wondering if someone could let me know if I made a dumb mistake?
0002 | 
0003 | Assume G is a finite group, let x,y be 2 elements in G, both having order 2.  Prove that <x,y> is either abelian or it is isomorphic to a dihedral group.  
0004 | 
0005 | My attempt:  Suppose <x,y> is not abelian.  We can imbed <x,y> into S\_n, the symmetric group on n letters, for some n >= 1.  So <x,y> is isomorphic to a subgroup H of  S\_n.  Under this isomorphism, x should get mapped to an element of order 2 and same for y.  So x and y are both mapped to transpositions (a b) and (c d).  As sets, We cannot have {a,b} and {c,d} be disjoint, as this would mean that the transpositions (a b) and (c d) would commute, so <x,y> would be abelian.  We cannot have (a,b) = (c, d) since we are imbedding <x,y>, , so WLOG by relabeling we can say (a b) = (1 2) and (c d) = (2 3).  Then H is isomorphic to S\_3, which is isomorphic to the dihedral group of order 6, so we are done.  
0006 | 
0007 | &#x200B;
0008 | 
0009 | This looks weird compared to some solutions I'm seeing online, was hoping someone could check if I made some error in this proof.  Thanks!
0010 | 
0011 | &#x200B;
\end{Vsrc}
\textbf{A YES means:} the proof verification document faithfully represents both the source context and the challenged claim, without dropping a condition or changing its meaning.
\section*{Question 2 --- Does the proposed counterexample disprove that claim?}
\textbf{Claim to test (repeated for convenience).}
\begin{Vsrc}
So x and y are both mapped to transpositions (a b) and (c d).
\end{Vsrc}
\subsection*{How the record says the argument uses this claim}
This is the record's description, not an established fact; check it against the source.
\begin{Vsrc}
Under an embedding into a symmetric group, every element of order 2 must map to a single transposition.
\end{Vsrc}
\subsection*{Proposed counterexample (exact values or construction)}
\begin{Vsrc}
Let G=C2={e,x}, with x^2=e. Define φ:G→S4 by φ(e)=id and φ(x)=(1 2)(3 4). Since φ(x)^2=id, φ is a homomorphism; since φ(x)≠id, it is injective. Thus φ is an embedding and x has order 2, but its image is two disjoint transpositions, not a single transposition.
\end{Vsrc}
\subsection*{Why the record says the counterexample works}
Check this explanation independently; it is evidence to assess, not an answer key.
\begin{Vsrc}
An order-2 permutation may be any nonempty product of disjoint transpositions. The explicit image (1 2)(3 4) has order 2 and cycle type (2,2), disproving the universal claim.
\end{Vsrc}
\subsection*{How to decide}
\begin{enumerate}
\item List every assumption that applies to the challenged claim.
\item Check that the proposed example satisfies every one of those assumptions.
\item Recompute the claimed outcome and the actual outcome using exact reasoning.
\item Confirm that the conflict is with the challenged claim itself, not an unrelated sentence.
\end{enumerate}
\textbf{A YES means:} the example is fully specified, satisfies all relevant assumptions, and makes the challenged claim false. If a necessary value, condition, or calculation is missing, choose CANNOT DETERMINE.
\section*{Question 3 --- Does the final Lean theorem statement verify this same counterexample?}
\textbf{What to decide.} A Lean file can compile while proving the wrong theorem. Compare these three links independently:
\begin{enumerate}
\item The Lean definitions represent the same mathematical objects, assumptions, and operations as the challenged claim.
\item The Lean construction uses the same concrete values or example shown in Question 2.
\item The final theorem statement asserts that the challenged claim fails on that example, not merely a weaker or nearby statement.
\end{enumerate}
\subsection*{Selected final Lean theorem statement (verbatim UTF-8 excerpt)}
This is an exact contiguous slice of the certificate file, from the theorem keyword through the original \texttt{:=} token where its proof begins. It is shown first so the formal target is not buried in the complete file.
\begin{Vlean}
theorem phi_x_is_not_single_transposition :
    Not (IsSingleTransposition (phi .x)) :=
\end{Vlean}
\subsection*{Project-local imported Lean files (complete, verbatim UTF-8)}
This section follows the certificate's import commands recursively and prints every repository-owned Lean source file on which the final theorem statement depends. Each file is copied in full, including its definitions and controls. External Mathlib, Std, and Lean Init source is not duplicated.
No project-local imported Lean files were found. External Mathlib, Std, and Lean Init source is not duplicated in this proof verification document.
\subsection*{Complete counterexample Lean file (verbatim UTF-8)}
The complete Lean file is included below so you can check every definition and supporting result used by the final theorem statement. Use it together with the project-local imported Lean files above.
\textbf{Lean certificate file:}
\begin{Vsrc}
partB_feasibility/reddit_certificates/reddit_learnmath_14s0k9c.lean
\end{Vsrc}
\begin{Vlean}
/-!
Record: reddit_learnmath_14s0k9c
Public source: https://www.reddit.com/r/learnmath/comments/14s0k9c/checking_for_error_in_group_theory_proof/

Exact target scope: the material local inference that, under an embedding into
a symmetric group, every element of exact order two must map to one single
transposition.  This certificate does not claim that the surrounding main
theorem is false: its nonabelian/two-generator context is not used by that
local inference, so it is deliberately not imposed on the concrete source
group.

Concrete independently verified witness: C2 = {e,x}, with x*x=e and x!=e,
embedded in S4 by phi(e)=id and phi(x)=(p0 p1)(p2 p3).  The image of x has
exact order two but is not one transposition.  The meaningful negative control
is the injective homomorphism psi:C2 -> S2 taking x to (q0 q1), which really is
a single transposition.

Material modelling judgments: C2 and both acted-on sets are explicit finite
inductive types.  C2 multiplication and inversion are full tables, and its
identity, associativity, and inverse laws are proved below.  `Perm` is a genuine
bijection representation containing forward and inverse maps with proofs of
both inverse composites.  `permComp p q` means p after q (q is applied first);
homomorphism preservation is stated pointwise, the extensional equality
appropriate to permutations.  Exact order two means nonidentity plus square
equal to identity.  `IsSingleTransposition` has its literal semantic meaning:
two distinct points are swapped and every other point is fixed.  No support
count or numeric cycle-type proxy is used.  The refutation states finiteness,
all C2 group laws, exact order of x, preservation, injectivity, exact order of
its image, and failure of the claimed conclusion.  The S2 control checks the
same substantive premises and the opposite conclusion.
-/

namespace RedditLearnmath14s0k9c

inductive C2 where
  | e | x
  deriving DecidableEq

inductive P4 where
  | p0 | p1 | p2 | p3
  deriving DecidableEq

inductive P2 where
  | q0 | q1
  deriving DecidableEq

def c2Mul : C2 -> C2 -> C2
  | .e, g => g
  | g, .e => g
  | .x, .x => .e

def c2Inv : C2 -> C2
  | .e => .e
  | .x => .x

def C2Exhaustive : Prop :=
  forall g : C2, g = .e \/ g = .x

def C2GroupLaws : Prop :=
  (forall g, c2Mul .e g = g /\ c2Mul g .e = g) /\
  (forall a b c, c2Mul (c2Mul a b) c = c2Mul a (c2Mul b c)) /\
  (forall g, c2Mul (c2Inv g) g = .e /\ c2Mul g (c2Inv g) = .e)

def ExactOrderTwoC2 (g : C2) : Prop :=
  g ≠ .e /\ c2Mul g g = .e

theorem c2_exhaustive : C2Exhaustive := by
  intro g
  cases g with
  | e => exact Or.inl rfl
  | x => exact Or.inr rfl

theorem c2_group_laws : C2GroupLaws := by
  refine And.intro ?_ (And.intro ?_ ?_)
  · intro g
    cases g <;> exact And.intro rfl rfl
  · intro a b c
    cases a <;> cases b <;> cases c <;> rfl
  · intro g
    cases g <;> exact And.intro rfl rfl

theorem x_has_exact_order_two : ExactOrderTwoC2 .x := by
  exact And.intro (by decide) rfl

structure Perm (alpha : Type) where
  toFun : alpha -> alpha
  invFun : alpha -> alpha
  leftInv : forall a, invFun (toFun a) = a
  rightInv : forall a, toFun (invFun a) = a

def permId {alpha : Type} : Perm alpha where
  toFun := fun a => a
  invFun := fun a => a
  leftInv := by intro a; rfl
  rightInv := by intro a; rfl

def permInv {alpha : Type} (p : Perm alpha) : Perm alpha where
  toFun := p.invFun
  invFun := p.toFun
  leftInv := p.rightInv
  rightInv := p.leftInv

/- Composition orientation: `permComp p q` applies q first, then p. -/
def permComp {alpha : Type} (p q : Perm alpha) : Perm alpha where
  toFun := fun a => p.toFun (q.toFun a)
  invFun := fun a => q.invFun (p.invFun a)
  leftInv := by
    intro a
    rw [p.leftInv, q.leftInv]
  rightInv := by
    intro a
    rw [q.rightInv, p.rightInv]

def PermEq {alpha : Type} (p q : Perm alpha) : Prop :=
  forall a, p.toFun a = q.toFun a

theorem perm_comp_assoc_pointwise {alpha : Type} (p q r : Perm alpha) :
    PermEq (permComp (permComp p q) r) (permComp p (permComp q r)) := by
  intro a
  rfl

theorem perm_id_left_pointwise {alpha : Type} (p : Perm alpha) :
    PermEq (permComp permId p) p := by
  intro a
  rfl

theorem perm_id_right_pointwise {alpha : Type} (p : Perm alpha) :
    PermEq (permComp p permId) p := by
  intro a
  rfl

theorem perm_inv_left_pointwise {alpha : Type} (p : Perm alpha) :
    PermEq (permComp (permInv p) p) permId := by
  intro a
  exact p.leftInv a

theorem perm_inv_right_pointwise {alpha : Type} (p : Perm alpha) :
    PermEq (permComp p (permInv p)) permId := by
  intro a
  exact p.rightInv a

def ExactOrderTwoPerm {alpha : Type} (p : Perm alpha) : Prop :=
  (exists a, p.toFun a ≠ a) /\ PermEq (permComp p p) permId

def IsSingleTransposition {alpha : Type} (p : Perm alpha) : Prop :=
  exists a b,
    a ≠ b /\
    p.toFun a = b /\
    p.toFun b = a /\
    forall z, z ≠ a -> z ≠ b -> p.toFun z = z

def swap4 : P4 -> P4
  | .p0 => .p1
  | .p1 => .p0
  | .p2 => .p3
  | .p3 => .p2

def doubleSwap : Perm P4 where
  toFun := swap4
  invFun := swap4
  leftInv := by intro a; cases a <;> rfl
  rightInv := by intro a; cases a <;> rfl

def phi : C2 -> Perm P4
  | .e => permId
  | .x => doubleSwap

def PreservesC2Mul {alpha : Type} (f : C2 -> Perm alpha) : Prop :=
  forall g h z,
    (f (c2Mul g h)).toFun z = (permComp (f g) (f h)).toFun z

def Injective {alpha beta : Type} (f : alpha -> beta) : Prop :=
  forall a b, f a = f b -> a = b

theorem phi_preserves_mul : PreservesC2Mul phi := by
  intro g h z
  cases g <;> cases h <;> cases z <;> rfl

theorem phi_injective : Injective phi := by
  intro a b h
  cases a <;> cases b
  · rfl
  · have bad : P4.p0 = P4.p1 :=
      congrArg (fun p : Perm P4 => p.toFun .p0) h
    cases bad
  · have bad : P4.p1 = P4.p0 :=
      congrArg (fun p : Perm P4 => p.toFun .p0) h
    cases bad
  · rfl

theorem phi_x_has_exact_order_two : ExactOrderTwoPerm (phi .x) := by
  refine And.intro ?_ ?_
  · exact Exists.intro P4.p0 (by decide)
  · intro a
    cases a <;> rfl

theorem phi_x_is_not_single_transposition :
    Not (IsSingleTransposition (phi .x)) := by
  intro h
  rcases h with ⟨a, b, hab, habForward, _hbaForward, hfix⟩
  have h0 := hfix P4.p0
  have h2 := hfix P4.p2
  cases a <;> cases b
  · exact hab rfl
  · have bad := h2 (by decide) (by decide)
    cases bad
  · cases habForward
  · cases habForward
  · have bad := h2 (by decide) (by decide)
    cases bad
  · exact hab rfl
  · cases habForward
  · cases habForward
  · cases habForward
  · cases habForward
  · exact hab rfl
  · have bad := h0 (by decide) (by decide)
    cases bad
  · cases habForward
  · cases habForward
  · have bad := h0 (by decide) (by decide)
    cases bad
  · exact hab rfl

abbrev PublishedClaimAtWitness : Prop :=
  C2Exhaustive ->
  C2GroupLaws ->
  ExactOrderTwoC2 .x ->
  PreservesC2Mul phi ->
  Injective phi ->
  ExactOrderTwoPerm (phi .x) ->
  IsSingleTransposition (phi .x)

/-! The source's local universal inference fails at the fully checked witness. -/
theorem reddit_learnmath_14s0k9c_refutation :
    C2Exhaustive /\
    C2GroupLaws /\
    ExactOrderTwoC2 .x /\
    PreservesC2Mul phi /\
    Injective phi /\
    ExactOrderTwoPerm (phi .x) /\
    Not (IsSingleTransposition (phi .x)) /\
    Not PublishedClaimAtWitness := by
  refine And.intro c2_exhaustive ?_
  refine And.intro c2_group_laws ?_
  refine And.intro x_has_exact_order_two ?_
  refine And.intro phi_preserves_mul ?_
  refine And.intro phi_injective ?_
  refine And.intro phi_x_has_exact_order_two ?_
  refine And.intro phi_x_is_not_single_transposition ?_
  intro claimed
  exact phi_x_is_not_single_transposition
    (claimed c2_exhaustive c2_group_laws x_has_exact_order_two
      phi_preserves_mul phi_injective phi_x_has_exact_order_two)

def swap2 : P2 -> P2
  | .q0 => .q1
  | .q1 => .q0

def singleSwap : Perm P2 where
  toFun := swap2
  invFun := swap2
  leftInv := by intro a; cases a <;> rfl
  rightInv := by intro a; cases a <;> rfl

def psi : C2 -> Perm P2
  | .e => permId
  | .x => singleSwap

theorem psi_preserves_mul : PreservesC2Mul psi := by
  intro g h z
  cases g <;> cases h <;> cases z <;> rfl

theorem psi_injective : Injective psi := by
  intro a b h
  cases a <;> cases b
  · rfl
  · have bad : P2.q0 = P2.q1 :=
      congrArg (fun p : Perm P2 => p.toFun .q0) h
    cases bad
  · have bad : P2.q1 = P2.q0 :=
      congrArg (fun p : Perm P2 => p.toFun .q0) h
    cases bad
  · rfl

theorem psi_x_has_exact_order_two : ExactOrderTwoPerm (psi .x) := by
  refine And.intro ?_ ?_
  · exact Exists.intro P2.q0 (by decide)
  · intro a
    cases a <;> rfl

theorem psi_x_is_single_transposition : IsSingleTransposition (psi .x) := by
  refine Exists.intro P2.q0 (Exists.intro P2.q1 ?_)
  refine And.intro (by decide) ?_
  refine And.intro rfl ?_
  refine And.intro rfl ?_
  intro z hz0 hz1
  cases z with
  | q0 => exact (hz0 rfl).elim
  | q1 => exact (hz1 rfl).elim

/-! Negative control: the same nonvacuous premises can have the claimed result. -/
theorem reddit_learnmath_14s0k9c_negative_control :
    C2Exhaustive /\
    C2GroupLaws /\
    ExactOrderTwoC2 .x /\
    PreservesC2Mul psi /\
    Injective psi /\
    ExactOrderTwoPerm (psi .x) /\
    IsSingleTransposition (psi .x) := by
  exact And.intro c2_exhaustive
    (And.intro c2_group_laws
      (And.intro x_has_exact_order_two
        (And.intro psi_preserves_mul
          (And.intro psi_injective
            (And.intro psi_x_has_exact_order_two
              psi_x_is_single_transposition)))))

end RedditLearnmath14s0k9c

\end{Vlean}
\subsection*{Proof-safety information}
The automated scan found no unfinished proofs, new axiom declarations, unsafe declarations, external code hooks, or uses of Lean's compiled decision procedure. This scan concerns proof construction only; it does not establish that the theorem matches the source claim.
\textit{Technical scan counts:} native\_decide=0, sorry=0, axiom\_decl=0, unsafe=0, extern=0. \textit{Technical tier:} KERNEL-ONLY.
\textbf{A YES means:} Lean represents the same objects and assumptions, uses the same proposed counterexample, and the final theorem statement asserts the same contradiction. Successful compilation by itself is not enough.
\clearpage
\section*{Human sign-off}
Complete this block from your own review of Questions 1--3. Do not use the automated appendix below as an answer key.\par
\textbf{Reviewer name:} \rule{9cm}{0.4pt}\par
\textbf{Date:} \rule{4cm}{0.4pt}\par
\medskip\textbf{1. Does the proof verification document accurately represent the source and challenged claim?}\par
\answerchoices\par
\decisionexplanation
\medskip\textbf{2. Does the proposed counterexample satisfy the assumptions and make that claim false?}\par
\answerchoices\par
\decisionexplanation
\medskip\textbf{3. Does the final Lean theorem statement represent and verify that same counterexample?}\par
\answerchoices\par
\decisionexplanation
\medskip\textbf{Overall decision:}\par
\checkbox\ VALIDATED (all three answers are YES)\par
\checkbox\ REJECTED (at least one answer is NO)\par
\checkbox\ NEEDS ANOTHER REVIEWER (at least one CANNOT DETERMINE and no NO)\par
\medskip\textbf{Additional comments.} Use this space for any extra source, counterexample, calculation, or Lean details that support your decisions.\par\rule{\linewidth}{0.4pt}\par\vspace{7pt}\rule{\linewidth}{0.4pt}
\clearpage
\section*{Appendix --- Prior automated review (background only)}
\textbf{These results are not an answer key.} They come from the earlier Opus/Fable verification campaign and are printed only so an auditor can see what was previously checked. The human decisions above must be made independently. Raw status codes are retained in small type for traceability.
\textbf{Overall automated disposition:} Cleared after a follow-up reviewer resolved an automated concern.
\par\textit{Audit code:} \nolinkurl{PASS_ADJUDICATED}
\subsection*{Independent search with no access to the stored example}
\textbf{Result:} The checker could not test this claim with a finite counterexample search.
\par\textit{Audit code:} \nolinkurl{CLAIM_NOT_FINITELY_REFUTABLE}
\par\textbf{Checker's explanation:}
\begin{Vsrc}
The blind record's problem field is only the title string 'checking for error in group theory proof' and supplies no mathematical content. The quoted claim is an intermediate proof step whose premises are entirely absent: the record names no group containing x and y, no relations or orders for x and y (it never says |x|=|y|=2), no map ('are mapped' has no named homomorphism, no source, no target, no injectivity hypothesis), no degree n for a symmetric group, and no meaning for the points a, b, c, d (nor whether {a,b} and {c,d} must be. (Shortened here; the complete note remains in the machine-check JSON.)
\end{Vsrc}
\par\textbf{Cases checked:}
\begin{Vsrc}
Diagnostic sweep only (no witness claimed): the full symmetric group S_n for n = 1,2,3,4,5,6 -- 873 permutations total, enumerated in itertools.permutations lexicographic order, of which 113 are involutions. Two schemas tested exhaustively over that space: S(n) 'every involution of S_n is a transposition' (all 113 involutions tested; (Shortened here; the complete note remains in the machine-check JSON.)
\end{Vsrc}
\subsection*{Independent re-check of the proposed counterexample}
\textbf{Result:} A separate checker recomputed the proposed counterexample and confirmed it.
\par\textit{Audit code:} \nolinkurl{CONFIRMED}
\par\textbf{Checker's note:}
\begin{Vsrc}
The witness is fully specified and finitary, and my independent recomputation reproduces every asserted fact. phi(x) = (1 2)(3 4) squares to the identity, so phi is a genuine homomorphism on C2; phi(x)!= id, so phi is injective and hence an embedding of a group containing an element of order 2 into S_4. Yet phi(x) has cycle type (2,2), not (2,1,1), so it is a product of two disjoint transpositions and not a single transposition. The universal claim therefore fails on the witness. (Shortened here; the complete note remains in the machine-check JSON.)
\end{Vsrc}
\subsection*{Independent review of the source-to-Lean correspondence}
\textbf{Result:} The Lean proof establishes a related failure, but at least one part of the correspondence needs review.
\par\textit{Audit code:} \nolinkurl{PARTIAL}
\par\textbf{Theorem summary from that reviewer:}
\begin{Vsrc}
`reddit_learnmath_14s0k9c_refutation` is an axiom-free conjunction over the explicit witness C2={e,x} -> Perm P4 with phi e = permId and phi x = doubleSwap = (p0 p1)(p2 p3): C2 is exhausted by {e,x}; the full C2 multiplication/inversion tables satisfy identity, associativity and inverse laws; x is nonidentity with x*x=e (exact order two); phi is a pointwise homomorphism (PreservesC2Mul phi) and injective; (Shortened here; the complete note remains in the machine-check JSON.)
\end{Vsrc}
\par\textbf{Difference noted by that reviewer:}
\begin{Vsrc}
Link 1 (definitions): clean. Perm is a genuine bijection record (toFun, invFun, both inverse laws), composition/identity/inverse/associativity are proved pointwise, ExactOrderTwoPerm = (has a moved point) AND (square is pointwise identity) which is exactly 'order exactly 2', and IsSingleTransposition carries its literal semantic content (two distinct points swapped, all others fixed) with no support-count or cycle-type proxy that could smuggle the conclusion. (Shortened here; the complete note remains in the machine-check JSON.)
\end{Vsrc}
\subsection*{Follow-up review of an automated concern}
\textbf{Result:} Follow-up review found that the record was right and an earlier checker was mistaken.
\par\textit{Audit code:} \nolinkurl{CHECKER_ERROR}.\par\textit{Confidence:} \nolinkurl{high}.\par\textit{Recommended action:} \nolinkurl{NONE}.
\par\textbf{Conclusion of the follow-up review:}
\begin{Vsrc}
The stored record is correct. In the source post x,y are elements of order 2 in a finite group with <x,y> nonabelian, <x,y> is embedded in some S_n, and the author writes 'So x and y are both mapped to transpositions (a b) and (c d).' Order-2 elements map to involutions, and an involution of S_n is any nonempty product of disjoint transpositions, not necessarily one. The record's witness is valid: phi: C2 -> S_4 with phi(x) = (1 2)(3 4) is a homomorphism (all four products checked), injective, phi(x) has order exactly 2 and cycle type (2,2), hence is not a single transposition. (Shortened here; the complete note remains in the machine-check JSON.)
\end{Vsrc}
\textit{This appendix is only background. You do not need any separate machine-output files to answer the three questions.}
\end{document}